\scanpage{91}
% archive foliation: 90
les \deleted{\unsure{objets}} épinglages de Legendre-Jacobi
``directs'', i.e. tels que la section de $\unsure{\boldsymbol{\mu}|\mu}_4^*$ qu'
\unsure{ils} définit soit justement $i$. \deleted{\unsure{Les}}
Le \inserted{champ des} courbes elliptiques relatives
\inserted{(sur des $S/\widetilde S_0$)} munies
d'un épinglage de Legendre-Jacobi direct,
est représentable par $U_{03\widetilde S_0}$, sur
lequel on fait opérer le groupe modulaire
$\mathbb{D}'_3$, qui est \inserted{justement} le sous-groupe
de $\mathbb{D}_3\times\unsure{\boldsymbol{\mu}|\mu}_{4\widetilde S_0}$ qui fait passer d'un
épinglage direct à un épinglage direct.

Revenant au cas de la base absolue
$S_0$, on doit encore expliciter les
opérations de $\unsure{\boldsymbol{\mu}|\mu}_4\times\unsure{\mathfrak{S}|\mathcal G}_3$ sur
$\underline{O}_{U_{0,3}\times\unsure{\boldsymbol{\mu}|\mu}_4^*}\deleted{\unsure{(E)}}(1)$,
via l'action tautologique résultant de
l'action de $\unsure{\mathfrak{S}|\mathcal G}_3$ sur
$V_{\{0,1,\infty\}}(\underline{O}_{S_0})$ d'
une part, de \inserted{$\unsure{\mathfrak{S}|\mathcal G}_3$ et} $\unsure{\boldsymbol{\mu}|\mu}_4$ sur le facteur $\unsure{\boldsymbol{\mu}|\mu}_4^*$ d'
autre part. On désigne l'action tautologique
d'un $(g,\lambda)$ \unsure{sur|par}
$\xi\in\unsure{\Gamma(\underline{O}_{U_{0,3}}(1)(A))}$
